\chapter{The political economy of engineered life}
\label{vol12:ch13}

\epigraph{An engineer is a political actor whether they wish it or not; the only choice is to act with awareness or without.}{Author}

\chapmeta{Half-life: medium-life. Archetypes: ethics, scaling. Exercise target: 15. Project track: research (read one major infrastructure decision's full record, produce a 5-page synthesis). Hazard class: medium (professional-conduct stakes; civilizational responsibility). Reader path default: standard, with the political-economy and the engineering-as-political sections core. Named cases: Robert Moses-era projects (\cite{hist:caro-powerbroker-1974}, registered Ch 7); the engineering profession's response to climate change over $30$ years (\cite{paper:franta-engineering-climate-2018}); the procurement-and-politics dynamics of major defence and civil projects (\cite{paper:flyvbjerg-megaprojects-2014}).}

\noindent
The chapter is the project's openly normative discussion of the engineering profession's political-economic context. The previous chapters of this volume have repeatedly arrived at the conclusion that the engineering profession's technical work is incomplete without engagement with the political-economic context that determines whether the technical work serves civilization or harms it. The chapter takes that conclusion as the topic in full, and is more openly argumentative than the previous chapters have been. The reader is asked to engage substantively with the chapter's arguments rather than to accept them; the chapter's purpose is to make the case for engineer-as-political-actor explicit so the reader can evaluate and respond to it.

\section{Who decides what gets built}\pathtag{core}

Engineering decisions are made by an interaction of: the engineering profession (the technical analysis; the design); the engineering employer (the firm or agency that commissions the work); the political-economic context (the regulatory regime; the financing structure; the political process); the affected populations (the users; the displaced; the downstream-impacted). The contemporary distribution of decision-making power across these actors varies by jurisdiction, sector, and project; the engineering profession is typically subordinate to the engineering employer in the immediate decision chain and typically embedded in the political-economic context whose decision-making constraints determine the work scope.

The engineering profession's traditional self-understanding has emphasised the technical role and de-emphasised the political role: the engineer as the technical contributor to decisions made by others. The chapter's argument: this self-understanding is incomplete at civilization scale, because the technical contributions matter politically and the engineer's choice to engage or not is itself a political choice. The chapter does not argue that engineers should become principally political actors; it argues that engineers should engage with the political-economic dimension of their work with the same discipline they bring to the technical dimension.

\begin{archetype}[Ethics: the engineer's degrees of freedom]
The engineer's degrees of freedom in a typical project: the technical-analysis-and-design choices within the scope; the engagement (or not) with the scope-definition process; the surfacing (or not) of consequences that lie outside the immediate scope; the engagement (or not) with the regulatory and standard-setting bodies; the engagement (or not) with the affected populations; the engagement (or not) with the engineering profession's collective bodies; the personal choice to accept or decline the work. The traditional engineering profession's training has emphasised the first; the chapter's argument is that the second through seventh degrees are equally part of the engineering scope at civilization scale.
\end{archetype}

\section{Who pays, who benefits, who is displaced}\pathtag{core}

Engineering projects' distributional consequences are typically well-known in principle and under-analysed in practice. A new highway: the users benefit; the displaced residents pay; the downstream pedestrians-and-cyclists pay; the future generations pay the climate-emission consequences. A new dam: the irrigation beneficiaries benefit; the displaced upstream communities pay; the downstream-fishery communities pay; the future-generations pay for the eventual reservoir-sedimentation problem. A new power plant: the customers benefit; the local-air-quality-impacted neighbourhood pays; the global-climate population pays the emissions consequence; the workers eventually displaced when the plant is retired pay the transition cost. The distributional pattern is the same across most major engineering projects, with the engineering decision determining the specific distribution.

The engineering profession's contribution to addressing the distributional consequences: the explicit distributional analysis as part of the project scope; the engagement with the affected populations in the scope-setting; the design choices that reduce the most negative distributional consequences; the engagement with the institutional structures (impact assessment; community-benefit agreements; just-transition support) that internalise the distributional consequences. The contemporary engineering practice across most jurisdictions includes substantial distributional analysis (environmental impact assessment; social impact assessment); the implementation varies in quality.

\begin{principle}[Distributional analysis as engineering scope]
A complete engineering scope at civilization scale includes the explicit distributional analysis: who pays, who benefits, who is displaced, who is downstream-impacted, who is intergenerational-impacted. The analysis is not an add-on to the technical scope but is part of the technical scope at civilization scale. The engineering profession's failure to integrate the distributional analysis is one of the consistent contemporary failure modes.
\end{principle}

\section{The maintenance economy and its political invisibility}\pathtag{standard}

Maintenance engineering is the unglamorous half of the engineering profession's work. The construction of new infrastructure is politically visible (the ribbon-cutting; the project announcement; the bond issue); the maintenance of existing infrastructure is politically invisible (the deferred replacement; the chronic-condition; the slowly-deteriorating-asset). The distribution of capital across the two: substantial under-investment in maintenance in most developed-world infrastructure, with the consequences documented in Vol X Chapter 14 (the Morandi Bridge; the Flint water crisis; the broader pattern). The engineering profession is differentially distributed across new-build and maintenance work, with the new-build side typically having more capital, more prestige, and more political support.

The contemporary engineering profession's contribution: the explicit asset-management discipline that quantifies the maintenance backlog, prioritises the work, and supports the political case for sustained maintenance funding. The professional advocacy for sustained maintenance investment (the ASCE Report Card; equivalent assessments in other countries) provides the engineering-credible documentation of the deficit. The institutional-political response has been mixed; some jurisdictions have substantially increased maintenance investment, others continue to defer.

\section{Standards and standards-setting bodies}\pathtag{mastery}

Standards (ASTM, ASME, IEEE, IEC, ISO, IETF, W3C, and a vast portfolio of more specific bodies) are the technical-codification infrastructure that makes engineering work interoperable. The standards-setting process is itself an engineering-and-political process: the technical contributions are made by the engineering profession; the decision-making is made by the institutional structures of the standards bodies; the participation requires the resources to engage. The under-representation of developing-country engineers and of public-interest perspectives in the major standards bodies is a chronic concern; the over-representation of industry-funded participants is the corresponding concern.

The engineering profession's responsibility includes the engagement with the standards-setting process at the level appropriate to the engineer's career stage and resources. The collective professional bodies (IEEE, ASME, ACM, ASCE, ICE, etc.) play a role in supporting the engagement; the contemporary expansion of public-interest technology fellowships and similar programmes is part of the response. The chapter's normative position: the standards-setting process is engineering-and-political at civilization scale, and the engineering profession bears responsibility for the quality of the participation.

\section{Procurement and how it shapes engineering}\pathtag{mastery}

Procurement is the contracting infrastructure through which engineering work is commissioned. The procurement regime determines: the scope-setting process (the requirements specification; the trade-off framework); the cost-evaluation method (low-bid; best-value; performance-based); the delivery model (design-bid-build; design-build; public-private-partnership); the risk-allocation between the engineering provider and the client; the documentation and accountability for the work product. The procurement choices substantially determine the engineering outcomes.

The Flyvbjerg megaproject literature \cite{paper:flyvbjerg-megaprojects-2014} documents the consistent pattern of cost-and-schedule overrun across major civil projects, with the median project exceeding budget by $\sim 30\%$-$50\%$ and schedule by $\sim 25\%$-$40\%$. The drivers: optimistic scope-setting at the procurement-design stage (the ``planning fallacy''); risk-allocation that produces incentive-misalignment; political-economic pressure for project approval that suppresses honest cost estimates; the contractor industry's adaptation to the bidding incentives. The engineering profession's contribution to addressing the pattern: the credible cost-and-schedule estimation methods; the participation in procurement-process reform; the engagement with the political-economic context that determines whether the procurement reform succeeds.

\section{The engineering profession and its institutional reality}\pathtag{standard}

The engineering profession's institutional reality is the network of: educational institutions (the engineering schools and accreditation bodies; ABET in the US; comparable elsewhere); professional licensing (the PE licence in the US; the comparable Chartered Engineer status in the UK and equivalent); professional societies (the discipline-specific societies, e.g., ASCE, ASME, IEEE; the cross-discipline societies, e.g., NSPE; the international federations, e.g., WFEO); regulatory bodies (the state engineering boards; the equivalent international bodies); employer organisations (the contracting firms; the consulting firms; the in-house engineering departments of major companies; the government engineering agencies); the engineering-press and continuing-education infrastructure. Each component is a separate institution with its own governance, incentives, and political position.

The contemporary engineering profession faces several institutional challenges: the demographic narrowness (the profession is substantially less diverse than the population it serves, with implications for the populations whose perspectives are represented in the engineering decisions); the workforce-supply constraints (the engineering-graduate-supply lag in many countries against demand growth); the credibility-and-relevance question (the perception of engineering as a technical-only profession that does not engage with the civilization-scale questions the contemporary moment requires); the climate-and-sustainability institutional response (the increasing professional-society engagement with climate; the unevenness across societies and countries). The engineering profession's collective response to these challenges over the next decades will determine whether the profession can credibly bear the civilization-scale responsibilities the project has documented.

\section{Engineering-as-political: the case made in full}\pathtag{core}

The chapter's central case: engineering at civilization scale is engineering, not policy, but the engineering-at-civilization-scale is also irreducibly political. The two propositions are compatible: the engineering profession brings technical discipline that is the engineering substance of civilization-scale work; the engineering profession is one of several actors whose decisions determine whether the technical substance is delivered, and the engineering profession's choice to engage or not engage with the political-economic context is itself a political choice with technical consequences.

The case is made positively (rather than negatively): engineers who engage with the political-economic context as part of their professional discipline do better engineering and produce better outcomes than engineers who do not. The cases that the previous chapters of this volume and of Volume X have documented are largely cases in which the engineering-political integration was inadequate, with consequences that the engineering profession has subsequently been asked to address. The cases that worked (the Montreal Protocol; the contemporary urban-water-utility asset management in well-managed jurisdictions; the contemporary aviation-safety regime under ICAO governance) are cases in which the engineering profession engaged credibly with the political-economic context and helped to shape the institutional structures that produced sustainable outcomes.

The case is also made with explicit modesty: the engineering profession is one of several actors and is not the principal political decision-maker. The chapter does not propose that engineers should determine civilizational direction; it proposes that engineers should engage with the political-economic context that determines whether their technical work delivers the civilizational benefit it could. The engineering profession's degrees of freedom are substantial but bounded; the chapter argues that the profession should use them with awareness rather than ignore them.

\begin{principle}[Engineer as political actor]
An engineer at civilization scale is a political actor whether they wish it or not. The technical contributions to the political-economic context determine the realisable outcomes; the engineering profession's choices about how and how much to engage shape the political-economic context's response. The engineering profession's traditional self-understanding as principally technical is inadequate to the civilization-scale responsibility the contemporary moment imposes. The discipline the chapter argues for: engaging with the political-economic dimension with the same disciplined approach the engineering profession brings to the technical dimension; recognising the boundaries of what engineering alone can deliver; not retreating into the technical-only role when the political-economic context demands engagement.
\end{principle}

\section{Failure: engineering decisions that hid the politics}\pathtag{core}

The chapter's failure cases each illustrate the consequences of the engineering profession's inadequate engagement with the political-economic context.

The Robert Moses-era urban renewal in New York (\cite{hist:caro-powerbroker-1974}; registered in Ch 7) is the case of engineering profession participating in projects whose political consequences were predictable but the engineers did not surface them. The Cross-Bronx Expressway, the Brooklyn-Battery Bridge proposal (defeated), the Westway controversy, and many other Moses-era projects involved engineering work whose technical execution was competent and whose distributional consequences were severe; the engineering profession's collective contribution at the time was to accept the technical scope without surfacing the distributional analysis.

The engineering profession's response to climate change over $30$ years \cite{paper:franta-engineering-climate-2018} is the case of slow profession-wide engagement: the climate-science findings were sufficiently established by the mid-1990s that the engineering profession could have engaged seriously with the decarbonisation engineering scope; the actual engagement across most engineering disciplines was modest until the 2010s; the contemporary substantial engagement is partly compensating for the earlier delay. The engineering profession's collective bodies (the professional societies; the engineering-education accreditation bodies) have been part of the slow-engagement pattern, with the contemporary acceleration partly driven by external pressure and partly by the profession's own re-engagement.

The Flyvbjerg megaproject overrun pattern \cite{paper:flyvbjerg-megaprojects-2014} is the case of profession-wide accommodation to political-economic pressure: the cost-and-schedule overrun pattern is substantially driven by the political-economic incentives at the procurement-and-approval stage, but the engineering profession has been part of the pattern by accepting the over-optimistic estimates rather than insisting on credible ones. The contemporary research-and-reform efforts (reference-class forecasting; project-cost benchmarking; the Major Projects Knowledge Hub at Oxford) are partly the profession's response.

\section{Project}\pathtag{core}

\begin{project}[Major-infrastructure-decision synthesis]
Choose a major infrastructure decision from the past $30$ years (a major transport project; a major energy project; a major water project; a major communications-infrastructure project; a major defence-acquisition decision). Read the full available record: the engineering documentation, the procurement record, the political-debate record, the downstream-effects documentation. Produce a $5$-page synthesis covering: the engineering scope and the engineering trade-offs as documented; the political-economic context of the decision-making; the distributional consequences as they emerged; the engineering profession's role in the decision-making and in the outcome; the lessons for engineering-political integration in similar future decisions. The project's purpose is to develop the reader's capacity for the engineering-political integration the chapter has argued for.
\end{project}

\section{Exercises}

\subsection*{Calculation}\pathtag{core}

\begin{exercise}
A major infrastructure project's cost estimate at the procurement stage is $\$2$ billion; the actual completed cost is $\$3.4$ billion. Compute the overrun fraction; compare to the Flyvbjerg-database median for the relevant project type.
\end{exercise}

\begin{exercise}
A maintenance backlog assessment for a US state's transportation infrastructure identifies $\$50$ billion of deferred work against an annual capital programme of $\$3$ billion. Compute the deferred-work-to-annual-capital ratio and identify the implications for asset condition.
\end{exercise}

\subsection*{Derivation}\pathtag{standard}

\begin{exercise}
Derive a distributional-analysis framework for a major infrastructure decision that identifies the present-and-future, on-site-and-off-site, direct-and-indirect cost-and-benefit categories. Identify the criteria for inclusion in the analysis.
\end{exercise}

\begin{exercise}
Derive a procurement-cost-and-schedule estimation framework based on reference-class forecasting. Identify the parameters and the data requirements.
\end{exercise}

\subsection*{Estimation}\pathtag{standard}

\begin{exercise}
Estimate the engineering-profession participation rate in standards-setting at major bodies (ASTM, ISO, IEEE). Identify the sectors and regions that are well-represented versus under-represented.
\end{exercise}

\begin{exercise}
Estimate the engineering-profession demographic distribution by gender, race-ethnicity, and country-of-origin in a major engineering discipline. Compare to the population the discipline serves and identify the implications.
\end{exercise}

\begin{exercise}
Estimate the engineering-profession's contribution to climate-policy formation over the past $30$ years (publications, professional-society positions, regulatory engagement). Identify the patterns of engagement and disengagement.
\end{exercise}

\subsection*{Design}\pathtag{standard}

\begin{exercise}
Design a distributional-impact assessment process for a major engineering organisation. Specify the scope, the methodology, the affected-population engagement, the documentation, the institutional accountability.
\end{exercise}

\begin{exercise}
Design a procurement-process reform for a public infrastructure agency facing chronic cost-and-schedule overruns. Specify the scope-setting, the estimation, the risk-allocation, the contract-management changes.
\end{exercise}

\begin{exercise}
Design a programme for sustained engineering-profession engagement with climate-and-sustainability decision-making at the institutional level. Specify the professional-society initiatives, the educational-curriculum changes, the regulatory-engagement infrastructure.
\end{exercise}

\subsection*{Judgement}\pathtag{mastery}

\begin{exercise}
Argue: the engineering profession's collective response to climate change over the past $30$ years has been inadequate; the profession bears responsibility for the slow engagement. Argue the opposite. Identify the institutional structures that determine the pace of professional engagement.
\end{exercise}

\begin{exercise}
Discuss the Flyvbjerg megaproject overrun pattern. Identify the engineering profession's role; identify the institutional changes that would reduce the pattern.
\end{exercise}

\begin{exercise}
Argue: engineer-as-political-actor is the appropriate professional self-understanding at civilization scale; the traditional technical-only self-understanding is inadequate. Argue the opposite. Identify the implications for engineering education and professional practice.
\end{exercise}

\begin{exercise}
Discuss the engineer's options when working on a project whose civilization-scale consequences the engineer believes are negative on net but the engineer's employer is committed to. Identify the institutional structures that protect or fail to protect the engineer's professional judgement.
\end{exercise}

\begin{exercise}
Discuss the engineering profession's demographic narrowness as a professional-practice and civilization-engineering concern. Identify the implications for the populations the profession serves; identify the institutional changes required.
\end{exercise}
